% GENERATED FILE. Edit canonical lessons, then run npm run generate:books.
% canonical-source-hash: fnv1a64:3236e694033d85e2
% canonical-lessons: RU-C03-byt, RU-C03-zhit, RU-C03-znat, RU-C03-govorit, RU-C03-videt, RU-C03-idti

\chapter{Six Verbs, and the One You Never Say}
\label{ch:russian-first-verbs}
\hlchaptermodality{3}

\begin{chapteropening}
\textbf{By the end of this chapter:} \emph{I can say what I do in Russian with six everyday verbs — \ru{быть}, \ru{жить}, \ru{знать}, \ru{говорить}, \ru{видеть}, \ru{идти} — put each of them in the \ru{я} form, negate one with \ru{не}, and say a sentence like «\ru{Я} — \ru{Анна}» the way Russian does, by leaving the verb out entirely.}

Say \ru{иду} — «I am going», one trip, right now — and hear why Russian splits going-there-once from going-there-often, a distinction English carries in tone and Russian carries in the verb itself. Then place \ru{идти} against the two verb families the chapter has been sorting, and meet the \ru{го}- that English also inherited in go / went.
\end{chapteropening}

\section[byt']{\ru{быть} — \textquotedblleft{}to be\textquotedblright{}}

\label{lesson:RU-C03-byt}



\begin{quote}
Chapter 2 ended on a Russian sentence with no verb in it at all. This lesson is the verb that wasn't there.
\end{quote}

\subsection*{You'll want to know first — \ru{быть}}
\begin{quote}
\textbf{\ru{быть}} — \emph{byt'} — \textbf{to be}
\end{quote}

Four letters, and you have met all four already: \textbf{\ru{б}} from \emph{\ru{спасибо}}, \textbf{\ru{ы}} from \emph{\ru{ты}}, \textbf{\ru{т}} from \emph{\ru{привет}}, and the soft sign \textbf{\ru{ь}} from \emph{\ru{очень}}. Nothing new to decode — which is rare, and worth enjoying.

The \textbf{\ru{ь}} is not a sound; it softens the \textbf{\ru{т}} in front of it. That ending \textbf{-\ru{ть}} is how almost every Russian verb names itself, the way English puts \emph{to} in front.

\begin{grammarlens}[title={Grammar Lens: the tense Russian will not say}]
English: \emph{I am a student.} Russian:

\begin{quote}
\textbf{\ru{Я} \ru{студент}.} — \emph{ya studént} — literally \textbf{\textquotedblleft{}I student.\textquotedblright{}}
\end{quote}

There is no word between them. Not a short one, not a quiet one — \textbf{none}. Russian does not use \emph{\ru{быть}} in the present tense; you link two things by setting them side by side.

The verb was not lost. Push the sentence into the past and it walks back in: \textbf{\ru{Я} \ru{был} \ru{студентом}} — \emph{ya byl studéntom} — \textquotedblleft{}I was a student.\textquotedblright{} The present is the one place it goes silent.

Chapter 1 already showed you the survivor: \textbf{\ru{нет}} is \emph{\ru{не}} + \emph{\ru{есть}}, \textquotedblleft{}not-is.\textquotedblright{} \emph{\ru{Есть}} means \textbf{there is}, never \emph{am}.
\end{grammarlens}

\begin{cousinweb}
\textbf{\ru{быть}} comes from Proto-Slavic *\emph{byti}, from PIE *\emph{b\textsuperscript{h}uH-}, \textquotedblleft{}to grow, to become.\textquotedblright{} Say \emph{byt'} and then say \textbf{be} — that is the same root, barely worn.

Its cousins: English \textbf{be} and \textbf{been}, German \emph{bin}, Latin \emph{fuī} \textquotedblleft{}I was\textquotedblright{} (which gave English \textbf{future}), Sanskrit \emph{bhávati}.

Now the part worth keeping. English \emph{be} is stitched from \textbf{three} different ancient roots — \emph{be} from *\emph{b\textsuperscript{h}uH-}, \emph{is} and \emph{am} from *\emph{h\textsubscript{1}es-}, \emph{was} and \emph{were} from *\emph{wes-}. Russian holds two of those apart and lets you watch: \textbf{\ru{быть}} is the *\emph{b\textsuperscript{h}uH-} one, \textbf{\ru{есть}} is the *\emph{h\textsubscript{1}es-} one — the same root as English \textbf{is}.
\end{cousinweb}

\subsection*{Your turn}
Say these aloud:

\begin{itemize}
\raggedright
  \item \textquotedblleft{}\ru{быть}\textquotedblright{} — \emph{byt'}, with a soft final t
  \item \textquotedblleft{}\ru{Я} \ru{студент}\textquotedblright{} — and feel the gap where English puts \emph{am}
  \item the two roots Russian keeps apart — \textquotedblleft{}\ru{быть} … \ru{есть}\textquotedblright{}
  \item the English pair — \textquotedblleft{}\textbf{be} … \textbf{is}\textquotedblright{}
\end{itemize}

\begin{tcolorbox}[breakable,colback=teal!4,colframe=teal!35!black,title={Before you move on}]
Say \textquotedblleft{}to be.\textquotedblright{} (\textbf{\ru{Быть}}.) Now say \textquotedblleft{}I am a student.\textquotedblright{} (\textbf{\ru{Я} \ru{студент}} — two words, and no verb.) Is the verb gone from the language? (No — \textbf{\ru{был}} returns in the past; only the present drops it.) Which English word is \emph{\ru{быть}}? (\textbf{Be} — same root, *\emph{b\textsuperscript{h}uH-}.) And which English word is \emph{\ru{есть}}? (\textbf{Is}.) Next: \ru{жить}, \textquotedblleft{}to live.\textquotedblright{}
\end{tcolorbox}
\section[zhit']{\ru{жить} — \textquotedblleft{}to live\textquotedblright{}}

\label{lesson:RU-C03-zhit}



\begin{quote}
Last lesson gave you a verb Russian refuses to say. This one you say constantly — and it is the first time you will put an ending on one.
\end{quote}

\subsection*{You'll want to know first — \ru{жить}}
\begin{quote}
\textbf{\ru{жить}} — \emph{zhit'} — \textbf{to live}
\end{quote}

Same \textbf{-\ru{ть}} ending as \emph{\ru{быть}}, and no new letters. The first one, \textbf{\ru{ж}}, you met buried in \emph{\ru{пожалуйста}}: it is the \emph{zh} of English \emph{mea\textbf{s}ure} or \emph{plea\textbf{s}ure} — a \emph{sh} with the voice switched on. English never gives that sound a letter of its own; Russian gives it one.

So: \emph{zhit'}, one syllable, soft \emph{t} at the end.

\begin{grammarlens}[title={Grammar Lens: the ending that means \textquotedblleft{}I\textquotedblright{}}]
\begin{quote}
\textbf{\ru{Я} \ru{живу}.} — \emph{ya zhivú} — \textbf{I live.}
\end{quote}

Two things happened. A \textbf{-\ru{у}} arrived on the end, and that \emph{-\ru{у}} is the whole signal for \textbf{I} — it is what \emph{am}, \emph{do} and \emph{-s} are in English, packed into one letter. And a \textbf{\ru{в}} surfaced in the middle that the infinitive never showed you: \emph{\ru{жи}-\ru{ть}} becomes \emph{\ru{жи}\textbf{\ru{в}}-\ru{у}}.

That is normal. A Russian verb's dictionary form and its working stem are allowed to differ, so you learn the pair, not the one.

Switch to \emph{you}, and the ending switches: \textbf{\ru{ты} \ru{живёшь}} — \emph{ty zhivyósh} — \textquotedblleft{}you live.\textquotedblright{} The person lives in the ending; Russian can therefore drop the pronoun entirely and still be perfectly clear.
\end{grammarlens}

\begin{cousinweb}
\textbf{\ru{жить}} goes back to PIE *\emph{g\textsuperscript{w}eih\textsubscript{3}-}, \textquotedblleft{}to live,\textquotedblright{} and the family is enormous: Latin \emph{vīvere} gave English \textbf{vivid, survive, revive, vital}; Greek \emph{bíos} gave \textbf{biology}, and Greek \emph{zō\'{\={e}}} gave \textbf{zoo} and \textbf{zodiac}.

But the cousin worth carrying is the plain English one: \textbf{quick}. It meant \textbf{alive}, not fast — the speed came later, from how the living move. It is still frozen in the phrase \emph{the \textbf{quick} and the dead}, in \emph{cut to the \textbf{quick}} (cut to the living flesh), and in \textbf{quicksilver}, the metal that behaves as if alive.

So \emph{\ru{жить}} is \textbf{quick} in its older sense: to be among the living.
\end{cousinweb}

\subsection*{Your turn}
Say these aloud:

\begin{itemize}
\raggedright
  \item \textquotedblleft{}\ru{жить}\textquotedblright{} — \emph{zhit'}, with the \emph{measure} sound in front
  \item \textquotedblleft{}\ru{Я} \ru{живу}\textquotedblright{} — and land on the \textbf{-\ru{у}}
  \item the pair — \textquotedblleft{}\ru{живу} … \ru{живёшь}\textquotedblright{}
  \item the old English meaning — \textquotedblleft{}\textbf{quick} … alive\textquotedblright{}
\end{itemize}

\begin{tcolorbox}[breakable,colback=teal!4,colframe=teal!35!black,title={Before you move on}]
Say \textquotedblleft{}to live.\textquotedblright{} (\textbf{\ru{Жить}}.) Say \textquotedblleft{}I live.\textquotedblright{} (\textbf{\ru{Я} \ru{живу}} — and the \textbf{-\ru{у}} is doing the work.) What did English \emph{quick} once mean? (\textbf{Alive} — as in \emph{the quick and the dead}.) Name two Latin cousins. (\emph{Vivid}, \emph{survive} — from \emph{vīvere}.) Next: \ru{знать}, \textquotedblleft{}to know.\textquotedblright{}
\end{tcolorbox}
\section[znat']{\ru{знать} — \textquotedblleft{}to know\textquotedblright{}}

\label{lesson:RU-C03-znat}



\begin{quote}
The most useful sentence a beginner owns is not \emph{hello}. It is \emph{I don't know} — and after this lesson you will have it.
\end{quote}

\subsection*{You'll want to know first — \ru{знать}}
\begin{quote}
\textbf{\ru{знать}} — \emph{znat'} — \textbf{to know}
\end{quote}

No new letters, and the same \textbf{-\ru{ть}} that ended \emph{\ru{быть}} and \emph{\ru{жить}}. Say the \textbf{\ru{з}} and \textbf{\ru{н}} together with no vowel between them — \emph{zn-}, exactly as English does in the middle of \emph{ma\textbf{gn}et} but never at the start of a word.

In the present it takes the ending you already own:

\begin{quote}
\textbf{\ru{Я} \ru{знаю}.} — \emph{ya znáyu} — \textbf{I know.}
\end{quote}

That \textbf{-\ru{ю}} is not a new ending. It is the same \textbf{-\ru{у}} you put on \emph{\ru{живу}}, spelled so it carries a \emph{y}-glide in front — the same glide \textbf{\ru{е}} carries in \emph{\ru{привет}}. Russian writes the glide into the vowel instead of adding a letter for it.

\begin{grammarlens}[title={Grammar Lens: \ru{не}, the whole of \textquotedblleft{}don't\textquotedblright{}}]
\begin{quote}
\textbf{\ru{Я} \ru{не} \ru{знаю}.} — \emph{ya ne znáyu} — \textbf{I don't know.}
\end{quote}

You met \textbf{\ru{не}} in Chapter 1, hidden inside \textbf{\ru{нет}} — \emph{\ru{не}} + \emph{\ru{есть}}, \textquotedblleft{}not-is.\textquotedblright{} Here it stands on its own, and it is the entire machinery of English \emph{don't}, \emph{doesn't} and \emph{didn't}: one small word, parked directly in front of the verb, never changing shape.

English needs a helper verb to negate — \emph{I know} becomes \emph{I \textbf{do} not know}, with a \emph{do} that appears from nowhere. Russian adds nothing and borrows nothing. Put \textbf{\ru{не}} in front, and you are done.
\end{grammarlens}

\begin{cousinweb}
\textbf{\ru{знать}} is PIE *\emph{ǵneh\textsubscript{3}-}, \textquotedblleft{}to know,\textquotedblright{} and so is English \textbf{know} — the silent \emph{k} in \emph{know} is the fossil of a sound Russian still pronounces in \emph{\ru{зн}-}.

Down the Latin branch: \emph{nōscere} gave \textbf{notice, notion, note, ignore} (\emph{not}-knowing), \textbf{recognize}, \textbf{incognito}, and — through \emph{nōbilis}, \textquotedblleft{}knowable, well-known\textquotedblright{} — \textbf{noble}. Down the Greek branch: \emph{gign\'{\={o}}skein} gave \textbf{diagnosis}, \textbf{prognosis}, \textbf{agnostic}. English \textbf{ken} and \textbf{cunning} come straight down the Germanic branch beside \emph{know}.
\end{cousinweb}

\subsection*{Your turn}
Say these aloud:

\begin{itemize}
\raggedright
  \item \textquotedblleft{}\ru{знать}\textquotedblright{} — \emph{znat'}, both consonants together
  \item \textquotedblleft{}\ru{Я} \ru{знаю}\textquotedblright{} — the \emph{-\ru{ю}} is your \emph{-\ru{у}} wearing a glide
  \item \textquotedblleft{}\ru{Я} \ru{не} \ru{знаю}\textquotedblright{} — and notice you added nothing else
  \item the family — \textquotedblleft{}\textbf{know} … \emph{notice} … \emph{ignore} … \emph{diagnosis}\textquotedblright{}
\end{itemize}

\begin{tcolorbox}[breakable,colback=teal!4,colframe=teal!35!black,title={Before you move on}]
Say \textquotedblleft{}to know.\textquotedblright{} (\textbf{\ru{Знать}}.) Say \textquotedblleft{}I don't know.\textquotedblright{} (\textbf{\ru{Я} \ru{не} \ru{знаю}}.) What does Russian add to negate a verb, besides \textbf{\ru{не}}? (\textbf{Nothing} — no helper verb at all.) Which English word is \emph{\ru{знать}}? (\textbf{Know} — and its silent \emph{k} is the sound Russian kept.) Next: \ru{говорить}, \textquotedblleft{}to speak.\textquotedblright{}
\end{tcolorbox}
\section[govorít']{\ru{говорить} — \textquotedblleft{}to speak\textquotedblright{}}

\label{lesson:RU-C03-govorit}



\begin{quote}
Your first new letter in three lessons, and an English word that looks like a relative and is not one.
\end{quote}

\subsection*{You'll want to know first — \ru{говорить}}
\begin{quote}
\textbf{\ru{говорить}} — \emph{govorít'} — \textbf{to speak, to talk}
\end{quote}

One letter here is \textbf{new}: \textbf{\ru{г}}. It is Greek \textbf{gamma} (Γ) with the corner rounded off, and it says a hard \textbf{g}, as in \emph{go} — never the \emph{j} of \emph{gem}.

Everything else you have met. And the two \textbf{\ru{о}}s behave the way \emph{\ru{спасибо}} taught you: an \textbf{\ru{о}} away from the stress slides towards \emph{a}. The stress falls on the last syllable, so this is \textbf{ga-va-RIT'}. Russian spells the \emph{o}; your mouth says \emph{a}.

\begin{grammarlens}[title={Grammar Lens: verbs come in two families}]
\begin{quote}
\textbf{\ru{Я} \ru{говорю}.} — \emph{ya gavaryú} — \textbf{I speak.}
\end{quote}

The \textbf{I} ending is the same \textbf{-\ru{ю}} you put on \emph{\ru{знаю}}. Every Russian verb in the present ends the same way for \emph{I}. Then switch to \emph{you}, and the two families separate:

\begin{itemize}
\raggedright
  \item \textbf{\ru{ты} \ru{знаешь}} — \emph{ty znáyesh} — you know
  \item \textbf{\ru{ты} \ru{говоришь}} — \emph{ty gavarísh} — you speak
\end{itemize}

An \textbf{-\ru{е}-} against an \textbf{-\ru{и}-}. That one vowel is the whole difference, and it is why you learn the \emph{you} form alongside the dictionary form: it names the family, and the remaining endings follow from it.
\end{grammarlens}

\begin{cousinweb}
\textbf{\ru{говорить}} is built on the old Slavic noun *\emph{govor\ru{ъ}}, \textquotedblleft{}a murmur, a din, the noise of many voices.\textquotedblright{} Russian still keeps \textbf{\ru{говор}} for the buzz of a crowd or a regional way of talking. Most etymologists read the root as \textbf{imitative} — the word sounds like what it names — which is honest but means there is no tidy English cousin to hand you. Where a root is uncertain, this book says so.

What there \emph{is} to hand you is a trap. \textbf{Govern} is \textbf{not} related. That is Latin \emph{gubernāre}, \textquotedblleft{}to steer a ship,\textquotedblright{} from Greek \emph{kybernân} — the same word that gave English \textbf{cybernetics} and every \textbf{cyber-} since. A governor steers; a \emph{\ru{говор}} is a noise. Two roots that never met.
\end{cousinweb}

\subsection*{Your turn}
Say these aloud:

\begin{itemize}
\raggedright
  \item the new letter — \textquotedblleft{}\textbf{\ru{г}} is gamma, and says \emph{g} as in \emph{go}\textquotedblright{}
  \item \textquotedblleft{}\ru{говорить}\textquotedblright{} — \emph{ga-va-RIT'}, both \emph{o}s leaning towards \emph{a}
  \item \textquotedblleft{}\ru{Я} \ru{говорю}\textquotedblright{} — same \textbf{-\ru{ю}} as \emph{\ru{знаю}}
  \item the two families — \textquotedblleft{}\ru{знаешь} … \ru{говоришь}\textquotedblright{}
\end{itemize}

\begin{tcolorbox}[breakable,colback=teal!4,colframe=teal!35!black,title={Before you move on}]
What sound is \textbf{\ru{г}}? (A hard \textbf{g}, from Greek gamma.) Say \textquotedblleft{}I speak.\textquotedblright{} (\textbf{\ru{Я} \ru{говорю}}.) Which vowel separates the two verb families in the \emph{you} form? (\textbf{-\ru{е}-} in \emph{\ru{знаешь}} against \textbf{-\ru{и}-} in \emph{\ru{говоришь}}.) Is \emph{\ru{говорить}} related to English \emph{govern}? (\textbf{No} — \emph{govern} is Greek \emph{kybernân}, \textquotedblleft{}to steer.\textquotedblright{}) Next: \ru{видеть}, \textquotedblleft{}to see.\textquotedblright{}
\end{tcolorbox}
\section[vídet']{\ru{видеть} — \textquotedblleft{}to see\textquotedblright{}}

\label{lesson:RU-C03-videt}



\begin{quote}
One root, two meanings, and half of them ended up in English as words about \textbf{knowing} rather than about eyes.
\end{quote}

\subsection*{You'll want to know first — \ru{видеть}}
\begin{quote}
\textbf{\ru{видеть}} — \emph{vídet'} — \textbf{to see}
\end{quote}

No new letters. Two you should keep honest, though: \textbf{\ru{в}} says \emph{v}, not \emph{b}, and \textbf{\ru{и}} says \emph{ee} — the pair that fooled you in \emph{\ru{привет}}.

Stress lands on the \textbf{first} syllable here — \emph{VEE-dyet'} — the opposite of \emph{\ru{говорит}}\ru{ь}. Russian never prints its stress, so it travels with the word.

The \emph{you} form tells you the family, exactly as promised last lesson:

\begin{quote}
\textbf{\ru{ты} \ru{видишь}} — \emph{ty vídish} — \textbf{you see}
\end{quote}

An \textbf{-\ru{ишь}}, so \emph{\ru{видеть}} belongs with \emph{\ru{говорить}}, not with \emph{\ru{знать}}.

\begin{grammarlens}[title={Grammar Lens: the letter that swaps for \textquotedblleft{}I\textquotedblright{}}]
Now the surprise:

\begin{quote}
\textbf{\ru{Я} \ru{вижу}.} — \emph{ya vízhu} — \textbf{I see.}
\end{quote}

The \textbf{\ru{д}} has become \textbf{\ru{ж}} — the \emph{measure} sound from \emph{\ru{жить}}. The ending is the ordinary \textbf{-\ru{у}}; it is the consonant in front of it that shifted.

This happens in the \textbf{I} form and nowhere else in the present. \emph{\ru{Ты} \ru{видишь}}, and everyone else, keep their \textbf{\ru{д}}. So the pattern to carry is small and exact: \textbf{\ru{вижу}}, then \textbf{\ru{видишь}} — and after that the verb behaves.

English does the same trick and has stopped noticing: \emph{speak} against \emph{speech}, \emph{break} against \emph{breach}, \emph{bake} against \emph{batch}. A consonant softening beside a front vowel, in both languages, for the same ancient reason.
\end{grammarlens}

\begin{cousinweb}
\textbf{\ru{видеть}} is PIE *\emph{weid-}, and that root carried \textbf{two} senses at once: \emph{to see}, and — because what you have seen is what you know — \emph{to know}.

The seeing branch, through Latin \emph{vidēre}: \textbf{video, vision, visit, evident, provide, review, advise}. Through Greek \emph{eîdos}: \textbf{idea} and \textbf{idol}.

The knowing branch came into English directly: \textbf{wit}, \textbf{wise}, \textbf{wisdom}, \textbf{witness} — someone who knows because they saw. Sanskrit took the same turn: the \textbf{Veda} is \textquotedblleft{}what has been seen,\textquotedblright{} meaning knowledge.

Russian kept both halves as two verbs: \textbf{\ru{видеть}} for the eyes, and the older \textbf{\ru{ведать}} for knowing — the \emph{ved-} of \emph{Veda}, still alive in Russian.
\end{cousinweb}

\subsection*{Your turn}
Say these aloud:

\begin{itemize}
\raggedright
  \item \textquotedblleft{}\ru{видеть}\textquotedblright{} — \emph{VEE-dyet'}, stress at the front
  \item the odd pair — \textquotedblleft{}\ru{вижу} … \ru{видишь}\textquotedblright{}
  \item which family it joins — \textquotedblleft{}\ru{видишь}, like \ru{говоришь}\textquotedblright{}
  \item the two branches — \textquotedblleft{}\emph{video} … \textbf{wit}\textquotedblright{}
\end{itemize}

\begin{tcolorbox}[breakable,colback=teal!4,colframe=teal!35!black,title={Before you move on}]
Say \textquotedblleft{}to see.\textquotedblright{} (\textbf{\ru{Видеть}}.) Say \textquotedblleft{}I see.\textquotedblright{} (\textbf{\ru{Я} \ru{вижу}} — the \textbf{\ru{д}} swapped to \textbf{\ru{ж}}.) Does that swap happen anywhere else in the present? (\textbf{No} — the \emph{I} form only.) What second meaning did *\emph{weid-} carry? (\textbf{To know} — hence \emph{wit}, \emph{wise}, \emph{Veda}.) Next: \ru{идти}, \textquotedblleft{}to go.\textquotedblright{}
\end{tcolorbox}
\section[idtí]{\ru{идти} — \textquotedblleft{}to go\textquotedblright{}}

\label{lesson:RU-C03-idti}



\begin{quote}
English has one verb for going. Russian has two, and picks between them on a distinction English hides inside its tenses.
\end{quote}

\subsection*{You'll want to know first — \ru{идти}}
\begin{quote}
\textbf{\ru{идти}} — \emph{idtí} — \textbf{to go} (on foot, and under way right now)
\end{quote}

No new letters, and only four of them. The \textbf{\ru{дт}} in the middle is not really pronounced as two sounds — Russians say \emph{it-TEE}, with the stress at the end.

The endings are ones you already have:

\begin{quote}
\textbf{\ru{Я} \ru{иду}.} — \emph{ya idú} — \textbf{I'm going.}
\end{quote}

>

\begin{quote}
\textbf{\ru{Ты} \ru{идёшь}.} — \emph{ty idyósh} — \textbf{you're going.}
\end{quote}

An \textbf{-\ru{у}} for \emph{I}, and an \textbf{-\ru{ешь}} ending for \emph{you}, so \emph{\ru{идти}} sits with \emph{\ru{знать}}, not with \emph{\ru{говорить}}.

\begin{grammarlens}[title={Grammar Lens: going now against going often}]
Russian keeps a \textbf{second} verb for the same activity:

\begin{quote}
\textbf{\ru{Я} \ru{иду}} — I am on my way, \textbf{now}, in one direction.
\end{quote}

>

\begin{quote}
\textbf{\ru{Я} \ru{хожу}} — \emph{ya khazhú} — I go, \textbf{as a habit}, back and forth.
\end{quote}

That is not a tense difference. They are two separate verbs, and Russian makes you choose one before you can open your mouth.

English draws the same line with grammar instead: \emph{I'm going} against \emph{I go}. You know the distinction perfectly — you have just never had to choose a \textbf{verb} for it. Every Russian verb of motion comes in such a pair, so this one teaches the pattern.
\end{grammarlens}

\begin{cousinweb}
\textbf{\ru{идти}} is PIE *\emph{h\textsubscript{1}ei-}, \textquotedblleft{}to go\textquotedblright{} — the same root as Latin \emph{īre}, which gave English \textbf{exit} (\textquotedblleft{}he goes out\textquotedblright{}), \textbf{transit}, \textbf{initial}, \textbf{itinerary}, \textbf{ambition} (\emph{ambīre}, \textquotedblleft{}to go around\textquotedblright{} canvassing votes) and \textbf{perish} (\emph{perīre}, \textquotedblleft{}to go through, to be lost\textquotedblright{}).

Now the good part. The past of \emph{\ru{идти}} is \textbf{\ru{шёл}} — \emph{shol} — which shares \textbf{no root} with it at all. A verb this common wears out and patches the hole with a second verb; that is called \textbf{suppletion}.

English does exactly this. \emph{Go} has no past of its own either: \textbf{went} was taken from \emph{wend} and kept. So \emph{\ru{идти} $\to$ \ru{шёл}} is \emph{go $\to$ went}, in two languages, for one reason — the words we use most wear out fastest.
\end{cousinweb}

\subsection*{Your turn}
Say these aloud:

\begin{itemize}
\raggedright
  \item \textquotedblleft{}\ru{идти}\textquotedblright{} — \emph{it-TEE}
  \item \textquotedblleft{}\ru{Я} \ru{иду}\textquotedblright{} — on my way this minute
  \item the contrast — \textquotedblleft{}\ru{иду} … \ru{хожу}\textquotedblright{}
  \item the two broken verbs — \textquotedblleft{}\ru{идти} … \ru{шёл}\textquotedblright{}, \textquotedblleft{}go … went\textquotedblright{}
\end{itemize}

\begin{tcolorbox}[breakable,colback=teal!4,colframe=teal!35!black,title={Before you move on}]
Say \textquotedblleft{}I'm on my way.\textquotedblright{} (\textbf{\ru{Я} \ru{иду}}.) And \textquotedblleft{}I go there regularly\textquotedblright{}? (\textbf{\ru{Я} \ru{хожу}} — a different verb, not a different tense.) What is the past of \emph{\ru{идти}}? (\textbf{\ru{Шёл}} — another root altogether, like English \emph{went}.) Name two English cousins of *\emph{h\textsubscript{1}ei-}. (\emph{Exit}, \emph{transit} — or \emph{itinerary}.) That is six Russian verbs, and one of them you will spend your life not saying.
\end{tcolorbox}
